\documentclass[aspectratio=169,11pt]{beamer}

% -------------------------------------------------
% Theta Join Kurzpraesentation (2 Folien)
% Rahmenbedingungen:
% - Keine Titelfolie
% - Eine Folie Erklaerung/Inhalte
% - Eine Folie Beispiel
% - Fusszeile: Name1, Name2, Name3 | Theta Join | Seitenzahl
% -------------------------------------------------

\usepackage[T1]{fontenc}
\usepackage[utf8]{inputenc}
\usepackage[ngerman]{babel}
\usepackage{lmodern}
\usepackage{amsmath,amssymb}
\usepackage{booktabs}

\input{env.tex}

\setbeamertemplate{navigation symbols}{}

% Eigene Fusszeile nach Vorgabe
\setbeamertemplate{footline}{%
  \leavevmode%
  \hbox{%
    \begin{beamercolorbox}[wd=.45\paperwidth,ht=2.5ex,dp=1ex,leftskip=1em]{author in head/foot}%
      \nameone, \nametwo, \namethree
    \end{beamercolorbox}%
    \begin{beamercolorbox}[wd=.35\paperwidth,ht=2.5ex,dp=1ex,center]{title in head/foot}%
      Theta Join
    \end{beamercolorbox}%
    \begin{beamercolorbox}[wd=.20\paperwidth,ht=2.5ex,dp=1ex,rightskip=1em plus1fil]{date in head/foot}%
      \insertframenumber/\inserttotalframenumber
    \end{beamercolorbox}%
  }%
  \vskip0pt%
}

\begin{document}

% Folie 1: Erklärung / Inhalte
\begin{frame}{Theta Join: Definition und Einordnung}
\textbf{Formale Definition:}

\begin{align*}
    R \bowtie_{\theta} S &:= \sigma_{\theta}(R \times S) \\
    [R \bowtie_{\theta} S] &:= [R] \cup [S]
\end{align*}

\textbf{Bedeutung:}
\begin{itemize}
    \item Zuerst Kreuzprodukt $R \times S$ bilden.
    \item Danach nach $\theta$ (Join Bedingung) filtern.
    \item $\theta$ kann Vergleiche wie $=$, $<$, $>$, $\leq$, $\geq$, $\neq$ und
    \item $\theta$ kann logische Verknuepfungen wie $R.a = S.b$ enthalten.
\end{itemize}

\end{frame}

% Folie 2: Beispiel
\begin{frame}{Beispiel zum Theta Join}
Gesucht: \textit{Vorname, Nachname} aller Mitarbeiter mit
Eintrittsdatum $> 20050201$ und Verguetungsstufe \texttt{Expert\_2}.

\vspace{0.3em}
\textbf{Relationale Algebra (mit Basisoperatoren):}
\begin{align*}
A &= Mitarbeiter \times Vergütungsgruppen \\
B &= \sigma_{\substack{Mitarbeiter.GehaltsStufe = Vergütungsgruppen.GehaltsStufe \\
     \land\ Eintrittsdatum > 20050201 \\
     \land\ Vergütungsgruppen.GehaltsStufe = 'Expert\_2'}}(A) \\
Ergebnis &= \Pi_{Nachname,\,Vorname}(B)
\end{align*}

\textbf{Mini-Beispieldaten:}
\begin{center}
\begin{tabular}{llll}
\toprule
Vorname & Nachname & Eintrittsdatum & GehaltsStufe \\
\midrule
Anna & Meier & 20060112 & Expert\_2 \\
Ben & Schulz & 20040210 & Expert\_2 \\
\bottomrule
\end{tabular}
\end{center}

\textbf{Ergebnis:} Nur \textit{Anna Meier} erfuellt alle Bedingungen.
\end{frame}

\end{document}
